% Generated draft: requires mathematical and source review.
\documentclass[11pt]{article}
\usepackage[margin=25mm]{geometry}
\usepackage{fontspec}
\setmainfont{Latin Modern Roman}
\usepackage{amsmath,amssymb,mathtools,unicode-math}
\setmathfont{Latin Modern Math}
\usepackage{enumitem,needspace,xurl,hyperref}
\setlength{\parindent}{0pt}
\setlength{\parskip}{5pt}
\setlength{\emergencystretch}{3em}
\newcommand{\N}{\mathbb{N}}
\newcommand{\Z}{\mathbb{Z}}
\newcommand{\Q}{\mathbb{Q}}
\newcommand{\R}{\mathbb{R}}
\newcommand{\C}{\mathbb{C}}
\newcommand{\F}{\mathbb{F}}
\newcommand{\A}{\mathbb{A}}
\newcommand{\PP}{\mathbb{P}}
\newcommand{\E}{\mathbb{E}}
\newcommand{\eps}{\varepsilon}
\newcommand{\dd}{\,\mathrm{d}}
\begin{document}
\section*{UnsolvedMath ID: 20000983}
\textbf{Source problem code:} catalogue entry
\subsection{Definitions and mathematical statement}
Caroline Klivans asks for a higher-dimensional analogue of Merino's graph identity relating the generating polynomial of critical configurations, graded by degree/level, to $T_G(1,y)$. David Perkinson asks about external activity. Klivans proposes the cellular matroid of a simplicial complex and an arithmetic Tutte refinement because higher critical-group orders involve weighted spanning-tree counts. Igor Pak proposes that degree might be constant on each firing coset, allowing a representative-free generating function. The source specifies neither the higher critical configurations nor a normalization of degree; even in the graph model the usual level grading involves a sink-dependent shift of total chip degree.

The aid supplies two obstructions to naive formulations. For any nonsingular integer reduced Laplacian $L\in\Z^{r\times r}$ and $w\in\Z^r$, the linear degree $c\mapsto w^{\mathsf T}c$ on $\Z^r$ descends to a well-defined integer-valued function on $\operatorname{coker}(L)=\Z^r/L\Z^r$ exactly when
\[
 w^{\mathsf T}L=0.
\]
Nonsingularity forces $w=0$. This obstructs a nonzero integer linear degree on the whole reduced firing quotient; it does not forbid a chosen-representative grading, a nonlinear statistic, or a degree with another target group.

For a reduced full-row-rank integral boundary matrix $M\in\Z^{r\times m}$ and $L=MM^{\mathsf T}$, let $r(A)$ be the rank of its columns indexed by $A$, and let $m(A)$ be the index of their lattice span in its saturation. The squared-multiplicity Tutte candidate is
\[
 T_M^{(2)}(x,y)=\sum_{A\subseteq\{1,\ldots,m\}}
 m(A)^2(x-1)^{r-r(A)}(y-1)^{|A|-r(A)}.
\]
At $(1,1)$, Cauchy--Binet gives
\[
 T_M^{(2)}(1,1)=\sum_{|A|=r}\det(M_A)^2=\det(MM^{\mathsf T})=|\operatorname{coker}(L)|,
\]
with zero determinants for dependent $A$. Thus this reduced-Gram model requires squared torsion/index weights for the natural count at $y=1$, rather than merely assuming the ordinary arithmetic multiplicity suffices. Root and lattice conventions can affect other higher-dimensional critical-group formulas.

The aid's explicit example is the boundary of a tetrahedron, rooted at edges $\{12,13,14\}$. It reports that Guzman--Klivans $(L,I)$ critical representatives, graded by raw coordinate sum, have enumerator $1+y+y^2+y^3$ in one orientation and $1+y^3+y^6+y^9$ after reversing retained edge $24$. In this example the ordinary, arithmetic, and squared-arithmetic Tutte polynomials agree and equal $x^3+x^2+x+y$, so all yield $3+y$ at $x=1$ and are unchanged by orientation. This claimed calculation rules out that naive raw-degree identity and the representative-free linear-degree shortcut, not every possible higher Merino theorem. The full orientation conventions and $(L,I)$ representative definition are absent and require review.
\subsection{Short English statement}
Find a higher-dimensional critical-configuration grading whose generating function extends Merino's Tutte identity. The supplied aid obstructs an integer linear grading on a reduced firing quotient and an orientation-dependent raw-degree proposal, while identifying squared lattice weights for the reduced-Gram counting normalization.
\subsection{Sources}
\begin{itemize}
\item AIM workshop \emph{Generalizations of chip-firing and the critical group}, problem 18, Caroline Klivans's question and comments by David Perkinson and Igor Pak: \url{https://aimath.org/pastworkshops/chipfiringproblems.pdf}. Merino is cited as reference [25], whose bibliographic entry is not supplied.
\item The attached aid names Guzman--Klivans $(L,I)$ critical configurations and the tetrahedral orientation example, without a full reference or complete representative conventions.
\item Local registry \texttt{data/unsolvedmath/problems.json}, permanent ID 20000983, exact source code \texttt{catalogue entry}; canonical location \texttt{aim-combinatorics-notes.json notes[107]}. Attached attempt 1, partial result, problem status at run partially solved, rigor proved, low novelty confidence; machine-generated, not peer reviewed, requiring independent verification before citation.
\end{itemize}
\end{document}
